% 1.6 — Planning
\section{Planning}
\label{sec:planning}

\subsection{Déroulement du stage}
Le stage a débuté le \datedebut{} et s'est terminé le \datefin.
\acompleter[phases du stage, de la prise en main de la bêta à la fin : à confirmer pour le
diagramme de Gantt d'ensemble]

\subsection{Chronologie des parties documentées}
L'historique du dépôt de travail permet de dater précisément les parties livrées entre le
21 septembre et le 9 octobre 2026. La figure~\ref{fig:gantt} les regroupe en cinq phases :
la reprise et l'audit ; un moteur qui prédit sans encore décider (ses prédictions sont
seulement écrites dans les reçus, pour être comparées aux décisions existantes) ; le moteur
qui décide ; l'extension famille par famille de la table de règles ; enfin la sécurité et
la sortie de l'ancien chemin par expressions régulières. La partie C3 (le moteur
PowerShell) est \statutCthree{}.

\begin{figure}[htbp]
\centering
% Figure 7 — Chronologie des parties livrées, d'après l'historique git (21/09 -> 09/10/2026).
% Une barre = les jours où les commits de la partie ont atterri.
\begin{ganttchart}[
    x unit=5.2mm, y unit chart=5.2mm, y unit title=5.5mm,
    time slot format=isodate, vgrid={*1{draw=trait,dotted}},
    title/.style={fill=grisclair,draw=trait},
    title label font=\sffamily\scriptsize,
    bar/.style={fill=wepwn!70,draw=wepwnfonce},
    bar height=0.55, bar label font=\sffamily\scriptsize,
    group/.style={fill=encre!80,draw=encre}, group label font=\sffamily\scriptsize\bfseries,
    group right shift=0, group top shift=0.55, group height=0.25, group peaks height=0.12,
    milestone/.style={fill=dASK,draw=dASK!70!black}, milestone label font=\sffamily\scriptsize,
    link/.style={-latex,draw=gris},
  ]{2026-09-21}{2026-10-11}
  \gantttitle{Jours du cycle de développement documenté}{21} \\
  \gantttitlelist{1,...,21}{1} \\
  \ganttgroup{Reprise et audit}{2026-09-21}{2026-09-27} \\
  \ganttbar{Import v1.7.0, outillage}{2026-09-21}{2026-09-22} \\
  \ganttbar{Audit : 18 défauts}{2026-09-26}{2026-09-27} \\
  \ganttbar{Modèle de commande}{2026-09-27}{2026-09-27} \\
  \ganttgroup{Moteur (prédiction seule)}{2026-09-28}{2026-09-29} \\
  \ganttbar{A, B, C1, C2a}{2026-09-28}{2026-09-28} \\
  \ganttbar{SYS1 : filets reliés}{2026-09-29}{2026-09-29} \\
  \ganttgroup{Le moteur décide}{2026-10-02}{2026-10-06} \\
  \ganttbar{S2+M1, M2.0+M2.1}{2026-10-02}{2026-10-02} \\
  \ganttbar{T, T5 : une politique}{2026-10-03}{2026-10-04} \\
  \ganttbar{B0, B1, B1b}{2026-10-05}{2026-10-06} \\
  \ganttbar{I1a--c, W0}{2026-10-06}{2026-10-06} \\
  \ganttgroup{Familles d'outils (M2.2)}{2026-10-06}{2026-10-08} \\
  \ganttbar{git, infra, cloud}{2026-10-06}{2026-10-07} \\
  \ganttbar{conteneurs, données, ASK}{2026-10-07}{2026-10-07} \\
  \ganttbar{envois, PaaS, SaaS}{2026-10-07}{2026-10-07} \\
  \ganttbar{système, LATER, préfiltre}{2026-10-08}{2026-10-08} \\
  \ganttgroup{Sécurité et sortie de l'ancien chemin}{2026-10-08}{2026-10-09} \\
  \ganttbar{W-SEC, V17}{2026-10-08}{2026-10-08} \\
  \ganttbar{M2.3a, M2.3b, V22, V18}{2026-10-09}{2026-10-09} \\
  \ganttbar[bar/.style={fill=white,draw=wepwnfonce,dashed}]{C3 (PowerShell)}{2026-10-10}{2026-10-11}
\end{ganttchart}

\caption{Chronologie des parties livrées, d'après l'historique git (21 septembre -- 9 octobre 2026)}
\label{fig:gantt}
\end{figure}

\begin{pourquoi}[Pourquoi un moteur « en ombre » avant de le laisser décider ?]
Les parties B à C2a ont livré un moteur qui \emph{ne décidait rien} : il écrivait sa
prédiction dans chaque reçu, et le portail vérifiait que zéro décision changeait. Pendant
ce temps, chaque divergence entre la prédiction et la décision existante était un
signal : soit un défaut du moteur, soit un défaut de l'ancien chemin. Le moteur n'a pris
la main (partie S2+M1) qu'une fois ces divergences comprises, et ce changement-là a été
revu ligne par ligne sur le portail.
\end{pourquoi}
